\documentclass[10pt,twocolumn]{article}
\usepackage[utf8]{inputenc}
\usepackage{geometry}
\usepackage{amsmath, amssymb}
\usepackage{graphicx}
\usepackage{booktabs}
\usepackage{enumitem}
\usepackage{xcolor}
\usepackage{mdframed}
\usepackage{multicol}
\usepackage{parskip}

\geometry{
  letterpaper,
  margin=0.45in,
  top=0.4in,
  bottom=0.4in
}

% Custom section: no numbering, blue bold, rule below
\makeatletter
\renewcommand{\section}{\@startsection{section}{1}{\z@}%
  {3pt \@plus 1pt \@minus 1pt}%
  {0.1pt}%
  {\small\bfseries\color{blue!70!black}}}
\renewcommand{\@seccntformat}[1]{}  % suppress section numbers
\makeatother

% Draw rule under section title using \addtocontents trick — simpler: wrap in a command
\newcommand{\mysection}[1]{%
  \section{#1}%
  \noindent\textcolor{blue!50!black}{\rule{\linewidth}{0.4pt}}\vspace{1pt}%
}

\setlength{\parskip}{1pt}
\setlength{\parindent}{0pt}
\setlength{\columnsep}{0.3in}

% Tighten align* vertical spacing globally
\setlength{\abovedisplayskip}{2pt}
\setlength{\belowdisplayskip}{2pt}
\setlength{\abovedisplayshortskip}{1pt}
\setlength{\belowdisplayshortskip}{1pt}

\newcommand{\E}{\vec{E}}
\newcommand{\B}{\vec{B}}
\newcommand{\F}{\vec{F}}
\newcommand{\dV}{\,dV}
\newcommand{\dl}{\,d\vec{l}}
\newcommand{\sn}[2]{#1\times10^{#2}}

% Compact equation environment
\newenvironment{eqbox}{%
  \begin{mdframed}[linewidth=0.5pt,innerleftmargin=3pt,innerrightmargin=3pt,
    innertopmargin=2pt,innerbottommargin=2pt,backgroundcolor=blue!4]
  \small
}{%
  \end{mdframed}\vspace{1pt}
}

\begin{document}
\footnotesize

{\centering\small\textbf{Physics Exam III --- Equation Sheet}\par
  \textit{Ch. 16--19 | Matter \& Interactions, Chabay}\par
  {ID: 0513544}\par}
\vspace{2pt}

%=============================================================
\mysection{Core Relationships: $\E$, $V$, and $I$}

\textbf{Electric Field} $\vec{E}$: force per unit charge; points from $+$ to $-$.\\
\textbf{Potential} $V$: energy per unit charge (scalar).\\
\textbf{Current} $I$: charge flow per unit time.

\begin{eqbox}
\textbf{Moving between them:}
\begin{align*}
  V_B - V_A &= -\int_A^B \vec{E}\cdot d\vec{l} \quad \text{(path independent)}\\
  \vec{E} &= -\nabla V \quad \text{(1-D shortcut: } E_x = -dV/dx\text{)}\\
  I &= \frac{dq}{dt} = nqAv_d \quad \text{\textbf{(Microscopic Current)}}\\
   &\quad \text{$n$: carrier density, $A$: area, $v_d$: drift speed}\\
  v_d &= \frac{I}{nqA} \quad \text{(drift speed; typically $\sim$mm/s)}\\
  \vec{J} &= nq\vec{v}_d \quad \text{(current density [A/m$^2$])}
\end{align*}
\end{eqbox}

\textbf{Key signs:} $\Delta V > 0$ moving against $\vec{E}$; $\Delta V < 0$ with $\vec{E}$.\\
Work by $\vec{E}$: $W = q\,\Delta V = q(V_i - V_f)$.

%=============================================================
\mysection{Ch. 16 -- Electric Potential \& Potential Energy}

\textit{Potential is a scalar; superpose by addition, not vector sum.}

\begin{eqbox}
\begin{align*}
  V &= \frac{1}{4\pi\epsilon_0}\frac{q}{r} \quad \text{(point charge)}\\
  U &= qV = \frac{1}{4\pi\epsilon_0}\frac{q_1 q_2}{r_{12}} \quad \text{(system PE)}\\
  \Delta V &= -E\,\Delta x \quad \text{(uniform field)}\\
  V &= -\int \vec{E}\cdot d\vec{l} \quad \text{(add contributions)}
\end{align*}
\end{eqbox}

``Electric Potential'' $\neq$ ``Electric Potential Energy'':\\
$U = qV$ (energy); $V$ alone is the potential (energy/charge).

%=============================================================
\mysection{Ch. 16 -- $V$ from $\E$ in Non-Uniform Fields}

\textit{Break path into pieces; $\vec{E}$ may vary along the path.}

\begin{eqbox}
\begin{align*}
  \Delta V &= -\int_i^f \vec{E}\cdot d\vec{l} \quad \text{(path independent)}\\
  V &= \frac{1}{4\pi\epsilon_0}\int \frac{dq}{r} \quad \text{(scalar, easier than } \vec{E}\text{)}
\end{align*}
Steps: \textbf{(1)} choose convenient path, \textbf{(2)} express $\vec{E}(l)$,\\
\textbf{(3)} evaluate $\int \vec{E}\cdot d\vec{l}$, \textbf{(4)} $\Delta V = -$result.
\end{eqbox}

%=============================================================
\mysection{Ch. 17 -- Electric Field in Circuits}

\textit{Mobile charges redistribute until steady state; the surface charge gradient creates $\vec{E}$ inside.}

\begin{eqbox}
\begin{align*}
  \vec{J} &= \sigma\vec{E} \quad \text{(Ohm, microscopic)}\\
  V &= IR, \quad R = \frac{\rho L}{A}, \quad \rho = \frac{1}{\sigma} \quad [\Omega\cdot\text{m}]\\
  v_d &= \frac{I}{nqA}, \quad |\vec{E}| = \frac{V}{L} \quad \text{(uniform)}
\end{align*}
\end{eqbox}

In steady state: $\vec{E}_{\perp\text{surface}}=0$ inside conductor; charge resides on surface.

%=============================================================
\mysection{Ch. 18 -- Circuits: Loop Rule \& Kirchhoff}

\textit{Closed loops conserve energy; junctions conserve charge.}

\begin{eqbox}
\begin{align*}
  \sum_{\text{loop}} \Delta V &= 0 \quad \textbf{(KVL -- Loop rule)}\\
  \sum I_{\text{in}} &= \sum I_{\text{out}} \quad \textbf{(KCL -- Junction rule)}\\
  R_{\text{eq}} &= \sum R_i \quad \text{(series)}\\
  \frac{1}{R_{\text{eq}}} &= \sum\frac{1}{R_i} \quad \text{(parallel)}\\
  P &= IV = I^2R = \frac{V^2}{R} \quad \text{[W]}
\end{align*}
\end{eqbox}

Sign convention: $-IR$ across resistor (in current direction); $+\mathcal{E}$ across EMF (negative to positive terminal).

%=============================================================
\mysection{Ch. 19 -- RC Circuits \& Capacitors}

\textit{Capacitors store charge; RC circuits decay/grow exponentially.}

\begin{eqbox}
\begin{align*}
  C &= \frac{Q}{V}, \quad C = \frac{\epsilon_0 A}{d} \quad \text{(parallel-plate)}\\
  U &= \tfrac{1}{2}CV^2 = \tfrac{Q^2}{2C} \quad \text{(energy stored)}\\
  Q(t) &= C\mathcal{E}\!\left(1-e^{-t/\tau}\right) \quad \textbf{(charging)}\\
  Q(t) &= Q_0\,e^{-t/\tau} \quad \textbf{(discharging)}\\
  \tau &= RC, \quad I(t) = \frac{\mathcal{E}}{R}e^{-t/\tau}
\end{align*}
Series: $\dfrac{1}{C_{\text{eq}}}=\sum\dfrac{1}{C_i}$ \quad Parallel: $C_{\text{eq}}=\sum C_i$
\end{eqbox}

%=============================================================
\mysection{Constants \& Quick Reference}

\begin{eqbox}
\begin{align*}
  \frac{1}{4\pi\epsilon_0} &= \sn{9.0}{9}\ \mathrm{N\cdot m^2/C^2}\\
  \epsilon_0 &= \sn{8.85}{-12}\ \mathrm{C^2/(N\cdot m^2)}\\
  e &= \sn{1.6}{-19}\ \mathrm{C}, \quad m_e=\sn{9.1}{-31}\ \mathrm{kg}
\end{align*}
\end{eqbox}

\begin{center}
\small
\begin{tabular}{ll}
\toprule
Quantity & SI Unit \\
\midrule
$V$ (potential) & V = J/C \\
$E$ (field) & N/C = V/m \\
$I$ (current) & A = C/s \\
$R$ (resistance) & $\Omega$ = V/A \\
$C$ (capacitance) & F = C/V \\
$\rho$ (resistivity) & $\Omega\cdot$m \\
$P$ (power) & W = J/s \\
\bottomrule
\end{tabular}
\end{center}

\vspace{2pt}
\textbf{Insulator vs.\ Conductor in $\vec{E}$:}\\
Conductor (static): $\vec{E}_{\text{inside}}=0$, $V=\text{const}$.\\
Insulator: $\vec{E}\neq 0$ inside; $\Delta V = \int \vec{E}\cdot d\vec{l} \neq 0$.
\begin{align*}
  E_x &= -\frac{dV}{dx} \quad \text{(gradient shortcut, 1-D)}
\end{align*}

\end{document}
